\documentclass[11pt]{amsart}
\usepackage[T1]{fontenc}
\usepackage{lmodern}
\usepackage{amsmath,amssymb,amsthm}
\usepackage{booktabs}
\usepackage[margin=1.1in]{geometry}
\usepackage[hidelinks]{hyperref}
\hypersetup{
  pdftitle={Log-convexity of strip dimer counts on the king's graph},
  pdfauthor={Dawid Ry\l{}ko},
  pdfsubject={MSC2020: Primary 05C70; Secondary 82B20, 15B48},
  pdfkeywords={king's graph, domicule tilings, perfect matchings, log-convexity, transfer matrix,
    Krein signature},
  pdfcreator={LaTeX with hyperref},
}
\usepackage{url}
\usepackage{microtype}
\emergencystretch=2em

\newtheorem{theorem}{Theorem}[section]
\newtheorem{lemma}[theorem]{Lemma}
\newtheorem{proposition}[theorem]{Proposition}
\newtheorem{conjecture}[theorem]{Conjecture}
\theoremstyle{definition}
\newtheorem{definition}[theorem]{Definition}
\newtheorem{problem}[theorem]{Open problem}
\theoremstyle{remark}
\newtheorem{remark}[theorem]{Remark}

\newcommand{\Z}{\mathbb{Z}}
\newcommand{\file}[1]{\texttt{#1}}

\title{Log-convexity of strip dimer counts on the king's graph}
\author{Dawid Ry\l{}ko}
\address{Independent researcher, Bielsko-Bia\l{}a, Poland}
\email{dawid.rylko@gmail.com}
\date{October 2026}
\thanks{This work is licensed under CC BY 4.0, \url{https://creativecommons.org/licenses/by/4.0/}.}
\subjclass[2020]{Primary 05C70; Secondary 82B20, 15B48}
\keywords{king's graph, domicule tilings, perfect matchings, log-convexity, transfer matrix,
  Krein signature}

\begin{document}

\begin{abstract}
Let $Z(k, m)$ be the number of perfect matchings of the $k \times m$ king board, the square grid
with both diagonals, i.e.\ the number of its domicule tilings. In all $1451$ cases we tested, on free
strips of width $m \le 26$ and cylinders of circumference $m \le 22$, the sequence $q \mapsto Z(2q, m)$
is log-convex; the tests are exact or carry a rigorous rounding bound. The property is not
structural. It fails on free strips of the triangular lattice, for every unequal pair of diagonal
weights we tried, where a Krein-signature argument for the left--right reflection predicts the
onset, and on wide strips for equal diagonal weights $x = \frac32, 2, 3$. Its limit, convexity of
$q \mapsto \log\lambda_{2q}$ for the strip growth rates, would give a conditional lower bound on the
entropy with ten correct digits. Along equal weights it is a competition between two modes of the
transfer matrix whose rates cross, on our data, between $x = \frac74$ and $x = 2$. We collect the obstructions to the
standard proof techniques and two exact reformulations.
\end{abstract}

\maketitle

\section{Introduction}

A \emph{domicule} is a pair of cells of the square grid that share a side or a corner, and a
domicule tiling of a board is a perfect matching of the \emph{king's graph} on its cells: cells are
adjacent when they are at chessboard king distance $1$. Let $Z(k, m)$ be the number of tilings of the
board with $k$ rows and $m$ columns (OEIS A239264 \cite{oeisA239264}), and $Z_{\mathrm{cyl}}(k, m)$
the same number when column $m - 1$ is also adjacent to column $0$. The companion paper
\cite{rylko2026entropy} computes the entropy per site $s = \lim \log Z(K, L)/(KL)$,
\begin{gather*}
  s = 0.5655119110413253 \pm 6\cdot10^{-16} \quad\text{(not rigorous)},\\
  0.545343749727 \le s \le 0.586216605404 \quad\text{(rigorous)}.
\end{gather*}
The rigorous bounds come from strips with free sides and converge only like $1/w$ in the width $w$.
On the square lattice, reflection positivity gives the Calkin--Wilf bounds
$(a_{2q+p} - a_{2q})/p \le s$ \cite{calkinwilf1998}, which converge exponentially; here
$a_k = \log\lambda_k$ and $\lambda_k$ is the growth rate of the free strip of width $k$,
$\lambda_k^2 = \lim_j Z(2j, k)^{1/j}$. Reflection positivity fails for the king's graph
(\S\ref{sec:obstructions}). The weakest hypothesis that still gives the even-$p$ bounds is
\begin{equation}\label{eq:Linf}
  \text{(L$_\infty$)}\qquad q \mapsto a_{2q} \text{ is convex}, \tag{L$_\infty$}
\end{equation}
since a convex sequence with $a_{2q}/(2q) \to s$ has increments that increase to $2s$. Under
(L$_\infty$), certified values of $\lambda_{22}$ and $\lambda_{24}$ give $s \ge 0.565511911035$
\cite[Prop.~8.2]{rylko2026entropy}. This note is about the finite statement behind it,
\begin{equation}\label{eq:L}
  \text{(L)}\qquad Z(2q, m)^2 \le Z(2q - 2, m)\, Z(2q + 2, m) \quad\text{for all } m \ge 1,\ q \ge 1. \tag{L}
\end{equation}
(L) implies (L$_\infty$): it makes $q \mapsto m^{-1}\log Z(2q, m)$ convex for every $m$; since
$Z(2q, m) = Z(m, 2q)$, these sequences tend to $a_{2q}$ along even $m$, and a pointwise limit of
convex sequences is convex.

We do \emph{not} conjecture (L). It holds in every case we tested at unit weights
(\S\ref{sec:tests}), but it is not a structural property of the model: it fails on the triangular
lattice, under asymmetric diagonal weights, and for larger equal diagonal weights
(\S\ref{sec:counter}). The conjecture of the companion paper is (L$_\infty$), and
\S\ref{sec:modes} explains why it is a quantitative statement about unit weights.

\section{Tests at unit weights}\label{sec:tests}

The counts $Z(k, m)$ were computed exactly with the row transfer of \cite{rylko2026entropy},
reduced modulo $60$-bit primes and combined by the Chinese remainder theorem; the number of primes
is chosen so that their product exceeds the Kahn--Lov\'asz bound on the number of perfect matchings,
so the result is the exact integer. For $m = 25, 26$ we used x87 extended precision ($64$-bit
mantissa) with a rigorous rounding bound: each cell step writes a sum of at most five nonnegative
terms, so the computed $Z(k, m)$ is within a factor $(1 + u)^{5mk}$ of the true value,
$u = 2^{-64}$, and the smallest certified margin is $1.237\cdot10^{-10}$. Table~\ref{tab:L} lists the
tests. (L) holds on free strips of width $m \le 26$ and on cylinders of circumference $m \le 22$ in
the stated row ranges. The margin $\log Z(2q-2) + \log Z(2q+2) - 2\log Z(2q)$ decreases monotonically
in $q$ for every $m$, by a ratio tending to $0.16$--$0.20$; there is no dip at intermediate $q$ and
nothing that shrinks with $m$.

\begin{table}[ht]
\caption{Tests of (L), $Z(2q)^2 \le Z(2q-2)Z(2q+2)$, at unit weights. The row range shrinks as the
width grows.}\label{tab:L}
\begin{tabular}{lllrr}
\toprule
lattice, boundary & widths $m$ & rows $2q$ & cases & violations \\
\midrule
king, free     & $2$--$24$  & $\le 120$ to $\le 20$ & 873 & 0 \\
king, cylinder & $3$--$22$  & $\le 80$ to $\le 22$  & 550 & 0 \\
king, free (x87, certified) & $25, 26$ & $\le 28$ & 28 & 0 \\
square, free / cylinder     & $\le 12$ & $\le 60$ & 319 / 290 & 0 / 0 \\
triangular, cylinder        & $\le 12$ & $\le 60$ & 290 & 0 \\
triangular, free            & $\le 12$ & $\le 60$ & 319 & \textbf{92} \\
\bottomrule
\end{tabular}
\end{table}

(L) is specific to even row counts: for even $m$, the sequence $j \mapsto Z(2j + 1, m)$ is
log-\emph{concave} in $217$ of the $220$ cases tested. This matches the free strips of the
companion paper, where even widths approach the limit of $a_k - sk$ from above and odd widths from
below.

\section{Counterexamples}\label{sec:counter}

\subsection{The triangular lattice}
On free strips of the triangular lattice (one diagonal per square) (L) fails from $m = 5$ on:
$Z(6, 5) = 6475$, $Z(8, 5) = 210560$, $Z(10, 5) = 6842465$ and
$210560^2 - 6475\cdot6842465 = 30552725 > 0$, confirmed by brute force. An argument for (L) must
therefore use the king's graph itself, not only non-planarity or non-bipartiteness.

\subsection{Unequal diagonal weights and the Krein signature}\label{sec:krein}
Give the orthogonal edges weight $1$ and the two diagonals weights $x_{\mathrm{dl}}$ and
$x_{\mathrm{dr}}$. The reflection $J\colon c \mapsto m - 1 - c$ of a strip of width $m$ exchanges the
two diagonals, so it is a symmetry exactly when $x_{\mathrm{dl}} = x_{\mathrm{dr}}$. Write the row
transfer as $T = HM$, with $H$ the horizontal dimers inside a row and $M$ the kernel of the edges
between two rows. Then $H = H^{\mathsf T} = JHJ$ and, for \emph{any} diagonal weights,
$M^{\mathsf T} = JMJ$: transposition exchanges the roles of the two rows, which exchanges the two
kinds of diagonals, and conjugation by $J$ exchanges them back. Hence $G = JM$ is symmetric and $GT = JMHM$ is symmetric, i.e.\ $T$
is self-adjoint for the indefinite form $G$, and $G e_0 = e_0$ for the empty state $e_0$.

\begin{lemma}\label{lem:krein}
Let $\lambda_j$ be a simple real eigenvalue of $T$ with right eigenvector $r_j$, and suppose
$r_j^{\mathsf T} G r_j \ne 0$. Then the weight of $\lambda_j$ in
$Z(k, m) = (T^k)_{00} = \sum_i c_i \lambda_i^k$ is
\[
  c_j = \frac{r_j(0)^2}{r_j^{\mathsf T} G\, r_j},
\]
so the sign of $c_j$ is the Krein signature of $r_j$.
\end{lemma}
\begin{proof}
From $GT = T^{\mathsf T} G$, the vector $G r_j$ is a left eigenvector for $\lambda_j$, so the spectral
projection is $P_j = r_j (G r_j)^{\mathsf T} / (r_j^{\mathsf T} G r_j)$, and
$c_j = (P_j)_{00} = r_j(0)\,(G r_j)(0) / (r_j^{\mathsf T} G r_j)$. Finally
$(G r_j)(0) = e_0^{\mathsf T} G r_j = (G e_0)^{\mathsf T} r_j = r_j(0)$.
\end{proof}

We checked the algebra in exact arithmetic for $m \le 7$ and five weight sets, and the formula
against a direct eigendecomposition for $m \le 9$ (relative error $7\cdot10^{-13}$). With equal
weights the empty state is $J$-even, $Z$ sees only the $J$-even sector, and there $G = M$; on the
$J$-odd sector $G = -M$. Under the perturbation $(1, 1 - \varepsilon)$ the top $J$-odd level, which
is $M$-positive, becomes visible with a \emph{negative} weight $-K\varepsilon^2$
($K = 0.70, 1.31, 1.63$ at $m = 4, 5, 8$, constant for $\varepsilon = 10^{-3}, \dots, 3\cdot10^{-2}$).

The sign of $Z(2q - 2)Z(2q + 2) - Z(2q)^2$ for large $q$ follows from the spectrum. Write
$Z(2q) = c_1 L_1^q + c_2 L_2^q + \dots$ with $L_i = \lambda_i^2$, where $L_1$ is the Perron root of
the irreducible part of $T^2$ that $e_0$ sees, and $L_2$ the visible level of next largest modulus.
If $L_2$ is real and positive and dominates the rest, then
\[
  Z(2q-2)Z(2q+2) - Z(2q)^2 \sim c_1 c_2 (L_1 L_2)^{q-1} (L_1 - L_2)^2 ,
\]
and $c_1 > 0$, so (L) fails for large $q$ exactly when $c_2 < 0$. The Perron vector is positive,
hence $J$-even; so a $J$-odd level can only matter if it lies above every visible level other than
the Perron root. Among $m \le 11$ this happens at $m = 4$ and $m = 8$; the onsets were computed for
$m \le 10$. The exact counts confirm the prediction: the weights $(1, \tfrac{9}{10})$ violate (L) at
$m = 4$ from $2q = 56$, at $m = 8$ from $2q = 86$ and at $m = 12$ from $2q = 112$, and
$(1, \tfrac{19}{20})$ at $m = 8$ from $2q = 118$. At $m = 4$ and $m = 8$ the onsets agree exactly with
the predicted ones; the failure at $m = 12$ lies outside the range of the prediction.

Clearly unequal weights fail much earlier. The weights $(1, \tfrac12)$ violate (L) at $m = 5$
($14$ of $19$ values of $q$, $2q \le 38$), $m = 8$ and $m = 9$, and at no other $m \le 9$;
$(1, \tfrac34)$ fails at $m = 9$ ($5$ of $19$ values, first at $2q = 30$). Every pair of unequal
diagonal weights on which (L) was tested, namely $(1, \tfrac{9}{10})$, $(1, \tfrac{19}{20})$,
$(1, \tfrac34)$, $(1, \tfrac12)$, $(1, \tfrac14)$, $(1, \tfrac{1}{10})$, $(1, 0)$ and
$(2, \tfrac12)$, violates it at some $m$. The symmetrised transfer $H(M + M^{\mathsf T})/2$ has no
violation where we tested it: for $m \le 9$, $2q \le 38$ and the weights $(1, \tfrac12)$,
$(1, \tfrac34)$, $(1, 0)$ and $(1, \tfrac{9}{10})$, and for $(1, \tfrac{9}{10})$ at $m = 8$ up to
$2q = 198$. This covers every failure of $HM$ in those ranges, but not the counterexamples at
$m = 4$ and $m = 12$ or the weights $(1, \tfrac{19}{20})$.

The orbits of $J$ split the count as $Z = Z_+ + Z_-$ with $Z_\pm = (Z \pm F)/2$, where $F$ counts the
$J$-symmetric tilings. $Z_+$ satisfies (L) for $m \le 12$; $Z_-$ does not (at $m = 4$ for every $q$,
otherwise only at the trivial $q = 1$, where $Z_-(0) = 0$).

\subsection{Equal diagonal weights}\label{sec:equal}
With equal weights $x$ on both diagonals, $J$ remains a symmetry, and (L) holds for
$x \in \{0, \tfrac14, \tfrac12, 1, 2, 4, 10, 100\}$ on narrow strips ($m \le 9$, $2q \le 38$), but
not on wide ones. Exactly, $Z(8, m)^2 > Z(6, m)\,Z(10, m)$ for every $m$ with $35 \le m \le 60$ at
$x = \tfrac32$, with $21 \le m \le 60$ at $x = 2$ and with $19 \le m \le 60$ at $x = 3$, while $x = 1$
and $x = \tfrac54$ give no violation at $2q = 8$ for $m \le 60$. These finite failures are exact
and do not depend on any extrapolation. So the left--right symmetry alone does not imply (L).

They are the finite-width traces of a failure of convexity of the growth rates at small $q$: for
$x \ge \tfrac32$ the second differences $a_{2q+2} - 2a_{2q} + a_{2q-2}$ are negative at $2q = 8$ and
$2q = 12$ (at $x = \tfrac32$: $-6.2\cdot10^{-4}$ at $2q = 8$). Whether (L$_\infty$), which is a
statement about large $q$, fails as well is the subject of the next section.

\section{Two modes and the crossing of their rates}\label{sec:modes}
At unit weights the second differences of $a_{2q}$ are positive for $2q = 2, \dots, 24$ ($a_0 = 0$),
those of the odd widths are negative for $k = 3, \dots, 23$, and the increments $a_{k+1} - a_k - s$
alternate in sign for $k = 0, \dots, 25$ \cite[\S8]{rylko2026entropy}. All of this is the
signature of one alternating mode, $a_k - sk \approx -2\sigma + C(-\rho)^k$ with $\rho = e^{-1.083}$
(the in-sector gap of the transfer) and $C > 0$. A second, faster mode comes from the splitting of
the two parity sectors of the free strip, which alternates with period $4$ in $k$ and decays by a
factor $-13.84$ per two columns. In the second differences
$d_q = a_{2q+2} - 2a_{2q} + a_{2q-2}$ the two modes appear as
\[
  d_q \approx P \rho_+^q + Q \rho_-^q, \qquad \rho_+ = \rho^2 = 0.1147,\quad \rho_- = -0.072,
\]
and at $2q = 24$ the second term is $0.007$ of the first, shrinking by a factor $0.63$ per step.
Along equal weights $x$ both rates grow, at different speeds. (L$_\infty$) holds for large $q$ when
$\rho_+ > |\rho_-|$ and $P > 0$, and fails when $|\rho_-| > \rho_+$. We define $x^*$ as the weight at
which the two rates cross, $\rho_+(x^*) = |\rho_-(x^*)|$.

\begin{table}[ht]
\centering
\caption{Two-mode Prony fits of the second differences $d_q$ of the free strips with both diagonal
weights $x$, widths $2q \le 22$. The ratio $|\rho_-|/\rho_+$ is given for the first and the last of
four sliding windows (window $2q = 8..14$ and $2q = 14..20$); the last column lists the even widths
$2q \le 20$ with $d_q < 0$.}\label{tab:modes}
\begin{tabular}{lllllll}
\toprule
$x$ & $\rho_+$ & $\rho_-$ & \multicolumn{2}{l}{$|\rho_-|/\rho_+$} & dominant & $d_q < 0$ at $2q =$\\
 & & & first & last & & \\
\midrule
$1$     & $0.1147$ & $-0.0718$ & $0.670$ & $0.627$ & positive & none\\
$5/4$   & $0.1433$ & $-0.1217$ & $0.873$ & $0.849$ & positive & none\\
$3/2$   & $0.1710$ & $-0.1611$ & $0.971$ & $0.942$ & positive & $8, 12, 16, 20$\\
$7/4$   & $0.1971$ & $-0.1938$ & $1.019$ & $0.984$ & positive & $8, 12, 16, 20$\\
$2$     & $0.2214$ & $-0.2219$ & $1.047$ & $1.002$ & alternating & $8, 12, 16, 20$\\
$3$     & $0.3027$ & $-0.3082$ & $1.095$ & $1.018$ & alternating & $8, 12, 16, 20$\\
\bottomrule
\end{tabular}
\end{table}

Table~\ref{tab:modes} gives the fits (\texttt{L/modes.py}, float64 power iteration to a relative
tolerance $10^{-15}$; the smallest $|d_q|$ used, $7.6\cdot10^{-10}$ at $x = 1$, is five orders of
magnitude above the error of the growth rates). At $x = 1$ the fitted $\rho_+ = 0.1147$ and
$\rho_- = -0.072$ agree with the two rates read off the unit-weight data above. Three features
matter.

\emph{At $x = 3/2$ the asymptotic sign is unresolved and leans convex.} The second differences are
negative at $2q = 8, 12, 16, 20$ and positive in between; the last negative value is
$-3.6\cdot10^{-9}$ at $2q = 20$. In the two-mode model the positive mode wins, $|\rho_-|/\rho_+ =
0.942$, and the model predicts $d_q > 0$ from about $2q = 24$ on. This is an extrapolation from four
windows, not a computation; what is exact at $x = 3/2$ are only the finite-strip failures of
Section~\ref{sec:equal}.

\emph{The ratio drifts downwards with the window} at every $x$, by $0.03$ to $0.08$ from the first
to the last window. At $x = 2$ and $x = 3$ the last ratios are $1.002$ and $1.018$, so the
alternating mode dominates on the data, and (L$_\infty$) would fail for large $q$ if these ratios
are the limits. That failure is plausible, but it is an extrapolation too: a further drift of the
same size would bring both ratios below $1$.

\emph{The crossing.} Linear interpolation of $\rho_+ - |\rho_-|$ between $x = 7/4$ and $x = 2$ gives
$x^* \approx 1.96$ from the last windows. Because of the drift this is a lower estimate of $x^*$
rather than a value; on the present data $x^* \in (7/4, 2)$, and widths beyond $22$ would be needed
to locate it better. Nothing in the data suggests a crossing at $x \le 3/2$.

The large-$x$ limit is suggestive. Diagonal steps preserve the colour of the checkerboard, so at
$x = \infty$ the model decouples into two square-lattice dimer models, one on each colour class
(each class with its diagonal edges is a square lattice turned by $45^\circ$), which are critical.
The alternating mode that takes over above $x^*$ may be a precursor of this decoupled limit. We do
not know whether $x^*$ marks a phase transition or only a crossover in the subleading spectrum.
Either way, $x = 1$ lies below $x^*$, and a proof of (L$_\infty$) at $x = 1$ has to see the size of the
diagonal weight; this rules out arguments that are uniform in $x$.

\section{Obstructions to standard proofs}\label{sec:obstructions}

\emph{Reflection positivity.} Reflect through a line between two rows. The reflection exchanges the
two diagonals, and the form that glues a top half to a bottom half is the number $M(A, B)$ of
matchings between the sets $A, B$ of cells matched across the line, with $a \sim b$ iff
$|a - b| \le 1$. This form is indefinite: for columns of length $m = 3$ its smallest eigenvalue is
$1 - \sqrt2$, with eigenvector $c = (1, -\sqrt2, 1)$ on $A = \{0\}, \{1\}, \{2\}$. Reflection
positivity would need $\langle F, \theta F\rangle \ge 0$ for every function $F$ of the top half. Take
$F = \sum_A (c_A / g(A))\,[\text{boundary state} = A]$, where $g(A) > 0$ is the number of three-row
half-boards with boundary state $A$ (here $10, 4, 10$); then $\langle F, \theta F\rangle = c^{\mathsf T}
M c < 0$. The vector $c$ is $J$-even, so restricting to the sector that $Z$ sees does not help: on the
$J$-even sector $M$ has $1, 1, 3, 8, 20, 34, 93, 186$ negative eigenvalues for $m = 3, \dots, 10$. By
the symmetry of the lattice the same holds for lines between columns. For a reflection through a line
of vertices the gluing kernel is $[U \ne D]$, which is indefinite (eigenvalues $\pm1$ already for
one vertex); we have not checked whether its negative directions are realised. At equal weights
$T = HM$ with $H$ and $M$ symmetric, so the obstruction is the indefiniteness of $M$, not the
asymmetry of $T$.

\emph{Moment sequences.} If $q \mapsto Z(2q, m)$ were a Stieltjes moment sequence of a positive
measure, (L) would follow; it is not. The eigenvalues of $T^2$ that $e_0$ sees are non-real or real
negative from $m = 3$ on (small: at most $0.047\,\lambda_1^2$ in modulus at $m = 10$), and the
Hankel matrices $[Z(2(i+j), m)]$ and $[Z(2(i+j+1), m)]$ have negative minors for every free width
$m \ge 3$ and every circumference $m \ge 7$. (L) is only the consecutive $2 \times 2$ part of this.

\emph{Total nonnegativity} of the transfer, in any ordering of the states, would force a real
nonnegative spectrum, and fails for the same reason. The
Lindstr\"om--Gessel--Viennot lemma needs a planar, acyclic path model that the crossing diagonals do
not provide. \emph{Interlacing} of the strip spectra would need a Hermitian structure, or a
Pontryagin space of small negative index, but $G$ has exponentially many negative directions.
Monotonicity of the gap gives the decay rate of the second differences, not their sign.

\section{Reformulations}\label{sec:reform}

\emph{Parity of paths.} Superpose a matching $X$ of rows $[0, 2q)$ with a matching $Y$ of rows
$[2, 2q+2)$ ($Z(2q)^2$ pairs), or $X'$ of $[0, 2q+2)$ with $Y'$ of $[2, 2q)$ ($Z(2q+2)Z(2q-2)$
pairs). The union $H$ has degree one on the top rows $\{0, 1\}$ and the bottom rows $\{2q, 2q+1\}$
and degree two in between, so it consists of paths and even cycles. In the first case top--top and
bottom--bottom paths have odd length and top--bottom (TB) paths even length; in the second all
paths are odd. Each $H$ has $2^{c(H)}$ labelings, $c(H)$ the number of cycles of length at least
four, so
\[
  Z(2q-2)Z(2q+2) - Z(2q)^2 = \sum_{H:\,t(H) \ge 1} 2^{c(H)}
  \bigl([\text{all TB paths odd}] - [\text{all TB paths even}]\bigr),
\]
where $t(H)$ is the number of TB paths. We checked this identity pair by pair for small $(m, q)$,
up to $3.4\cdot10^{6}$ pairs. The odd configurations win separately for every $t$ (e.g.\ $m = 5$,
$q = 2$: for $t = 2$, $1847660$ against $912656$; for $t = 4$, $32616$ against $0$). So (L) asks for
a sign-reversing injection that flips the parity of all top--bottom paths, and by
\S\ref{sec:counter} it would have to depend on the lattice and on the size of the diagonal weight.
With a shift of one row instead of two, the same identity computes $Z(2q+1)Z(2q-1) - Z(2q)^2$,
which is negative for even $m$; there the even configurations win.

\emph{A Krein-space form.} Cut the $2q$-row board in the middle and reflect the bottom half (the
lattice is invariant under $r \mapsto -r$). Then $Z(k_1 + k_2) = u_{k_1}^{\mathsf T} M u_{k_2}$,
with $u_k$ the boundary vector of a $k$-row half-board and $M$ the symmetric, indefinite gluing form
above. Since $Z(2q \pm 2) = u_{q \pm 1}^{\mathsf T} M u_{q \pm 1} > 0$, (L) says exactly that $M$ is
positive semidefinite on the span of $u_{q-1}$ and $u_{q+1}$: the Gram determinant of the pair is
$Z(2q-2)Z(2q+2) - Z(2q)^2$. The Hankel matrices of \S\ref{sec:obstructions} are the $M$-Gram matrices
of the longer families $\{u_0, u_2, u_4, \dots\}$ and $\{u_1, u_3, \dots\}$, on which $M$ is not
positive.

\emph{Partial involutions.} Swapping the parts of two superposed matchings below the first row
where their boundaries agree (plainly or after the reflection $J$) is an involution between the
two types of pairs, preserves the number of diagonal dimers, and accounts for part of the
difference. It leaves an unmatched remainder that does not shrink with $q$: the ratio of the two
unmatched counts tends to $1$ (from $0.26$ to $0.98$ at $m = 3$ and from $0.54$ to $0.99999$ at
$m = 5$, for $q = 2, \dots, 8$). Together with the flip of the pairs without top--bottom paths, it
reduces the even-type pairs with $t \ge 1$ only from $194$ to $104$ ($m = 3$) and from $912656$ to
$861364$ ($m = 5$). With asymmetric weights the twisted swap breaks the identity, as it must.

\section{Open problems}

\begin{problem}
Prove (L$_\infty$) at unit weights, or its consequence $a_{24} - a_{22} \le 2s$. By
\S\ref{sec:modes} this amounts to the sign $C > 0$ of the amplitude of the mode $C(-\rho)^k$ in
$a_k - sk$, with the faster alternating mode and the remainder controlled for every even $k \ge 22$.
\end{problem}

\begin{problem}
Locate $x^*$ more precisely, and decide whether (L$_\infty$) fails at $x = 2$ and $x = 3$, where the
two-mode fits favour the alternating mode, and holds at $x = \frac32$, where they favour the bulk
mode. Both need growth rates of wider strips than those used here.
\end{problem}

\section*{Reproducibility}
{\sloppy All scripts and outputs are in the repository \url{https://github.com/dawidrylko/lattice-dimers},
directory \texttt{research/king-entropy/}; every number of this note is mapped to its output file in
\file{paper/NUMBERS.md}. The tests of \S\ref{sec:tests} are \file{src/zexact.py},
\file{src/ltest.py}, \file{src/lcert.py} and \file{src/cw.py}; the weights of \S\ref{sec:counter}
are \file{src/lweight.py}, \file{L/symmetry.py} and \file{L/linf.py}; the two modes of
\S\ref{sec:modes} are \file{L/modes.py}; the obstructions and reformulations are
\file{src/lspec.py}, \file{src/lparity.py} and \file{L/involution.py}. The laboratory notes with the
history of these experiments are \file{notes/L-NOTE.md} and \file{notes/CW-NOTE.md}.\par}

\section*{Data availability statement}
The code and data that support the findings of this note are openly available in the repository
\url{https://github.com/dawidrylko/lattice-dimers}.

\section*{Disclosure statement}
The author reports no competing interests.

\bibliographystyle{amsplain}
\bibliography{refs}

\end{document}
